\chapter{関連研究}

\section{本章の概要}

本章では、本研究に関連する先行研究を、VLA と Physical AI、長期タスクと履歴利用、VLM による操作支援、Active Perception、Human-Robot Interaction、失敗検出、不確実性下の計画の観点から整理し、本研究の位置づけを明確にする。

\section{VLA と Physical AI}

RT-2 は web 規模知識をロボット制御へ転移する VLA の方向性を示し、OpenVLA はオープンソースでの実装可能性を広げた~\cite{rt2, openvla}。また、RT-1 は大規模実機データに基づく行動生成の基礎を築いている~\cite{rt1}。これらの研究は、視覚・言語・行動を統合し、ロボットに汎用的な操作能力を与える枠組みとして重要である。一方で、本研究は VLA そのものを訓練するのではなく、VLA や primitive 実行の前段で、現在工程の不確実性を診断する補助モジュールを対象とする。

\section{長期タスク・履歴・記憶を扱う研究}

TraceVLA、MemoryVLA、MAP-VLA は、長期ロボット操作において履歴、軌跡、記憶を利用する研究である~\cite{tracevla, memoryvla, mapvla}。これらは長期タスクにおける状態維持や文脈保持の重要性を示している。本研究も履歴と現在工程を扱う点では共通するが、目的は行動生成性能の向上ではなく、同じ不確実性が工程文脈によって異なる処理分岐へ割り当てられるかを診断することにある。

\section{VLM によるロボット操作支援}

MOKA は、VLM に画像上のマークを与えることで、自然言語指示から操作に必要なキーポイントや affordance を推定する~\cite{moka}。この方向は、VLM を直接操作支援へ利用する有効性を示している。しかし、本研究では VLM に直接操作点や次行動を出させるのではなく、工程条件が判定可能か、不確実性が何に由来するか、その不確実性をどの処理分岐へ送るべきかを診断する点が異なる。

\section{Active Perception と追加観測}

AP-VLM などの研究は、能動的に観測や探索を行うことで、操作に必要な情報を取得する枠組みを提案している~\cite{apvlm}。これらは、視覚証拠不足に対して reobserve を選ぶ発想に近い。一方で、本研究は追加観測だけでなく、ask\_user、stop、act、defer を含めて分類対象とするため、不確実性解消手段をより広い分岐集合として扱う。

\section{Ask-user と Human-Robot Interaction}

KnowNo は、LLM プランナが不確実な場合に人間へ助けを求める枠組みを示している~\cite{knowno}。また、GestLLM や GestOS などの HRI 研究は、人間由来情報をロボット制御へ接続する方向を発展させている~\cite{gestllm, gestos}。本研究は ask\_user のみを目的とせず、不確実性が人間確認に属するのか、安全停止に属するのか、あるいは後回し可能なのかを工程文脈に基づいて判定する点に特徴がある。

\section{失敗検出・実行監視・過程理解}

AHA は、VLM を用いてロボット操作の失敗を検出し、その理由を説明する研究である~\cite{aha}。RoboProcessBench は、VLM がロボット操作過程を理解できるかを評価するベンチマークを提供している~\cite{roboprocessbench}。これらが主として失敗後や過程理解を扱うのに対し、本研究は次工程へ進む前の工程条件不確実性を対象とする。

\section{不確実性下の計画と World Model}

TAMPURA は、不完全観測や行動結果の不確実性を含む長期計画を扱い、情報収集や不可逆リスクの回避を計画へ組み込んでいる~\cite{tampura}。また、V-JEPA 2 や FLARE は、将来予測や world model を通じて embodied AI の頑健性向上を目指している~\cite{vjepa2, flare}。本研究はこれらの大規模計画機構を直接提案するものではなく、現在の工程条件に関する不確実性の種類と処理分岐を VLM で診断する局所的かつ実務的な層を扱う。

\section{本研究の位置づけ}

本研究は、VLA、長期履歴利用、Active Perception、ask-user、failure detection、TAMP、world model の中間に位置する。既存研究が個別に扱ってきた「見る」「聞く」「止まる」「実行する」という機能を、長期タスク中の工程文脈依存的不確実性診断として統合的に整理し、同一視覚状態・異なる工程文脈における処理分岐の切り替えを評価対象とする点に独自性がある。
